% =====================================================================
%  config/polices.tex
%  Configuration typographique (fontspec — XeLaTeX).
%
%  Choix : polices 100 % LIBRES, déjà présentes dans TeX Live / la
%  distribution de polices Noto. Aucune police propriétaire.
%
%  Remarque sur la charte UNIGE : la police institutionnelle officielle
%  est « TheSans » (LucasFonts), qui est PROPRIÉTAIRE et ne peut pas être
%  redistribuée dans ce modèle. On utilise donc « Noto Sans » comme
%  substitut libre pour les éléments de type titre / logotype. Si vous
%  disposez d'une licence TheSans installée sur votre système, voir la
%  note « SUBSTITUTION THESANS » plus bas.
% =====================================================================

\usepackage{fontspec}
\defaultfontfeatures{Ligatures=TeX} % «--», «---», guillemets… gérés correctement

% ---------------------------------------------------------------------
%  1) POLICE PRINCIPALE (corps du texte)
%     Noto Serif : couverture Latin + Cyrillique + Grec (utile pour le
%     russe et le grec en écriture « inline » sans changer de police).
%
%  Repli automatique (voir la NOTE sur les polices variables plus bas et
%  scripts/get-fonts.sh) : si des faces STATIQUES de Noto Serif ont été
%  déposées dans ressources/polices/, on les charge par chemin ; sinon
%  on charge « Noto Serif » par son nom (cas normal sur Overleaf).
% ---------------------------------------------------------------------
\IfFileExists{ressources/polices/NotoSerif-Regular.ttf}{%
  \setmainfont{NotoSerif}[
    Path            = ressources/polices/,
    Extension       = .ttf,
    UprightFont     = *-Regular,
    BoldFont        = *-Bold,
    ItalicFont      = *-Italic,
    BoldItalicFont  = *-BoldItalic,
  ]%
}{%
  \setmainfont{Noto Serif}%
}

% ---------------------------------------------------------------------
%  2) POLICE SANS EMPATTEMENT (titres, page de titre, logotype)
%     Substitut libre de TheSans. Même repli statique que ci-dessus.
% ---------------------------------------------------------------------
\IfFileExists{ressources/polices/NotoSans-Regular.ttf}{%
  \setsansfont{NotoSans}[
    Path            = ressources/polices/,
    Extension       = .ttf,
    UprightFont     = *-Regular,
    BoldFont        = *-Bold,
    ItalicFont      = *-Italic,
    BoldItalicFont  = *-BoldItalic,
  ]%
}{%
  \setsansfont{Noto Sans}%
}

% ---------------------------------------------------------------------
%  3) POLICE À CHASSE FIXE (code, chemins, exemples techniques)
% ---------------------------------------------------------------------
\IfFileExists{ressources/polices/NotoSansMono-Regular.ttf}{%
  \setmonofont{NotoSansMono}[
    Path=ressources/polices/, Extension=.ttf,
    UprightFont=*-Regular, BoldFont=*-Bold,
    Scale=MatchLowercase,
  ]%
}{%
  \setmonofont{Noto Sans Mono}[Scale=MatchLowercase]%
}

% ---------------------------------------------------------------------
%  SUBSTITUTION THESANS (optionnel)
%  Si vous possédez TheSans, commentez le \setsansfont ci-dessus et
%  décommentez :
%
%  \setsansfont{TheSans}[
%    UprightFont = TheSans-Plain,
%    BoldFont    = TheSans-Bold,
%    ItalicFont  = TheSans-PlainItalic,
%  ]
% ---------------------------------------------------------------------

% =====================================================================
%  POLICES POUR LES ÉCRITURES NON LATINES
%  Déclarées comme familles séparées ; polyglossia (langues.tex) les
%  associe aux langues correspondantes. On garde aussi des commandes
%  \...font manuelles pour un usage direct.
% =====================================================================

% =====================================================================
%  NOTE IMPORTANTE SUR LES POLICES CJK ET ARABE ET LE PDF
%  ---------------------------------------------------------------------
%  XeLaTeX produit le PDF via « xdvipdfmx », qui NE SAIT PAS embarquer
%  une police VARIABLE (fichiers du type « NotoSerifCJK-VF.ttc » ou
%  « NotoNaskhArabic[wght].ttf »). Si votre système ne fournit QUE des
%  versions variables (cas de certaines distributions Linux récentes,
%  p. ex. Fedora), la compilation échoue avec :
%        xdvipdfmx:fatal: Invalid font
%
%  Sur OVERLEAF et la plupart des installations TeX Live, les versions
%  STATIQUES de Noto sont présentes : le chargement par NOM ci-dessous
%  fonctionne directement. Voir le README (section « Polices ») pour le
%  repli sur les systèmes à polices variables.
% =====================================================================

% --- Arabe (écriture de droite à gauche) -----------------------------
%  Noto Naskh Arabic : style naskh classique, très lisible.
%  Alternatives libres : « Amiri », « Scheherazade New ».
%
%  Repli automatique : si des instances STATIQUES ont été déposées dans
%  ressources/polices/ (voir scripts/get-fonts.sh et le README), on les
%  charge en priorité ; sinon on charge la police système par son nom.
\IfFileExists{ressources/polices/NotoNaskhArabic-Regular.ttf}{%
  \newfontfamily\arabicfont[
    Path        = ressources/polices/,
    Extension   = .ttf,
    UprightFont = NotoNaskhArabic-Regular,
    BoldFont    = NotoNaskhArabic-Bold,
    Script      = Arabic,
  ]{NotoNaskhArabic}%
}{%
  \newfontfamily\arabicfont{Noto Naskh Arabic}[Script=Arabic]%
}

% --- Chinois (simplifié) ---------------------------------------------
\IfFileExists{ressources/polices/NotoSerifCJKSC-Regular.otf}{%
  \newfontfamily\chinesefont[
    Path        = ressources/polices/,
    Extension   = .otf,
    UprightFont = NotoSerifCJKSC-Regular,
    BoldFont    = NotoSerifCJKSC-Bold,
  ]{NotoSerifCJKSC}%
}{%
  \newfontfamily\chinesefont{Noto Serif CJK SC}%
}

% --- Japonais --------------------------------------------------------
\IfFileExists{ressources/polices/NotoSerifCJKJP-Regular.otf}{%
  \newfontfamily\japanesefont[
    Path        = ressources/polices/,
    Extension   = .otf,
    UprightFont = NotoSerifCJKJP-Regular,
    BoldFont    = NotoSerifCJKJP-Bold,
  ]{NotoSerifCJKJP}%
}{%
  \newfontfamily\japanesefont{Noto Serif CJK JP}%
}

% --- Grec et Russe / Cyrillique --------------------------------------
%  Noto Serif couvre déjà le grec et le cyrillique ; on définit des
%  familles explicites pour que polyglossia dispose de \greekfont et
%  \russianfont dédiés. Même repli statique que la police principale.
\IfFileExists{ressources/polices/NotoSerif-Regular.ttf}{%
  \newfontfamily\greekfont[
    Path=ressources/polices/, Extension=.ttf,
    UprightFont=*-Regular, BoldFont=*-Bold,
    ItalicFont=*-Italic, BoldItalicFont=*-BoldItalic,
  ]{NotoSerif}%
  \newfontfamily\russianfont[
    Path=ressources/polices/, Extension=.ttf,
    UprightFont=*-Regular, BoldFont=*-Bold,
    ItalicFont=*-Italic, BoldItalicFont=*-BoldItalic,
  ]{NotoSerif}%
}{%
  \newfontfamily\greekfont{Noto Serif}%
  \newfontfamily\russianfont{Noto Serif}%
}
